% This is a complete, runnable LaTeX document.
\documentclass[11pt, a4paper]{article}

% --- PACKAGES ---
\usepackage[utf8]{inputenc}
\usepackage[T1]{fontenc}
\usepackage[margin=1in]{geometry} % Sets nice 1-inch margins
\usepackage{listings}             % The main package for typesetting code
\usepackage{xcolor}               % Required for defining custom colors
\usepackage{caption}              % Allows captions for listings
\usepackage{wrapfig}              % Allows figure wrapping

% --- CUSTOM COLOR DEFINITIONS ---
% We define some colors for our syntax highlighting
\definecolor{codegreen}{rgb}{0,0.6,0}
\definecolor{codeblue}{rgb}{0,0,0.6}
\definecolor{codegray}{rgb}{0.5,0.5,0.5}
\definecolor{backcolour}{rgb}{0.98,0.98,0.98}

% --- LISTINGS STYLE FOR ALLOY ---
% We define a new style called "alloy"
\lstdefinestyle{alloy}{
    backgroundcolor=\color{backcolour},   % Set a light gray background
    commentstyle=\color{codegreen},       % Set comment color
    keywordstyle=\color{codeblue},        % Set keyword color
    basicstyle=\ttfamily\small,          % Use a small, monospaced font
    breakatwhitespace=false,             % Don't break lines only at whitespace
    breaklines=true,                     % Automatically break long lines
    captionpos=b,                        % Put captions at the bottom
    frame=single,                        % Draw a single frame around the code
    keepspaces=true,                     % Keep spaces to preserve indentation
    numbers=left,                        % Add line numbers on the left
    numberstyle=\tiny\color{codegray},   % Style for line numbers
    showspaces=false,                    % Don't show spaces as '␣'
    showstringspaces=false,              % Don't show spaces in strings
    showtabs=false,                      % Don't show tabs

    % --- Alloy Specifics ---
    morecomment=[l]{//},                 % Define '//' as a single-line comment
    morecomment=[s]{/*}{*/},             % Define '/*' and '*/' as block comments
    morekeywords={
        abstract, sig, extends,    % Signatures
        one, lone, some, set,      % Multiplicities
        fact, all, no, in, not,    % Logic and facts
        pred, fun, run, check, for % Commands
    }
}

% --- DOCUMENT START ---
\begin{document}

\title{Typesetting Alloy Code in \LaTeX}
\author{A \LaTeX{} Example}
\date{}
\maketitle

Here is a simple example of an Alloy model. We use the \texttt{listings} package to format the code, which includes highlighting for both single-line (\texttt{//}) and block-style (\texttt{/* ... */}) comments, as well as keywords.

% We call the 'lstlisting' environment and tell it to use our 'alloy' style
\begin{lstlisting}[
    language=Alloy, % This is not a real language, but applies our style
    style=alloy,
    caption={A simple Alloy model of a file system.},
    label={lst:filesystem}
]
/*
 * This is a block comment.
 * This model describes a basic file system.
 */
abstract sig FSObject {} // An abstract filesystem object

// A File is a type of FSObject
sig File extends FSObject {}

// A Directory is also an FSObject
sig Dir extends FSObject {
    contents: set FSObject // A directory contains a set of FSObjects
}

// A fact to prevent a directory from containing itself (no loops)
fact NoSelfContainment {
    all d: Dir | d not in d.^contents
}

// Run the model to find an instance
run {} for 3
\end{lstlisting}

We can easily reference our code, for example, in Listing~\ref{lst:filesystem}. This shows how to define signatures like \texttt{File} and \texttt{Dir}, and how to write a \texttt{fact} that applies to the model.

\end{document}